\documentclass[10pt]{beamer}

\usepackage{appendixnumberbeamer}
\usepackage{booktabs}
\usepackage[scale=2]{ccicons}
\usepackage{pgfplots}
\usepackage{xspace}
\usepackage{bookmark}
\usepackage{amssymb}
\usepackage{mathtools}
\usepackage[normalem]{ulem}
\usepackage[T1]{fontenc}
\usepackage[sfdefault,book]{FiraSans}
\usepackage{FiraMono}
\usepackage{fontawesome}
\usepackage{listings}
\usepackage{hyperref}
\usepackage[10pt]{moresize}

\renewcommand*\familydefault{\ttdefault}

\renewcommand{\thefootnote}{\fnsymbol{footnote}}
\renewcommand{\MintedPygmentize}{/home/thiago/.local/bin/pygmentize}

\definecolor{light-gray}{gray}{0.95}
\renewcommand\lstlistingname{Código}
\DeclareCaptionFormat{listing} {
  \parbox{\textwidth}{\hspace{-0.2cm}#1#2#3}
}
\DeclareCaptionFont{black}{\color{black}}
\captionsetup[lstlisting]{
  format=listing,
  labelfont=black,
  textfont=black,
  singlelinecheck=true,
  margin=0pt,
  font={tt,footnotesize,bf}
}

\lstset{
  basicstyle=\footnotesize\ttfamily,
  escapeinside={\%*}{*)},
  mathescape=true,
  showspaces=false,
  showtabs=false,
  showstringspaces=false,%
  % backgroundcolor=\color{light-gray},
  rulesepcolor=\color{black},
  frame=shadowbox,
  literate=
  {á}{{\'a}}1 {é}{{\'e}}1 {í}{{\'i}}1 {ó}{{\'o}}1 {ú}{{\'u}}1
  {Á}{{\'A}}1 {É}{{\'E}}1 {Í}{{\'I}}1 {Ó}{{\'O}}1 {Ú}{{\'U}}1
  {à}{{\`a}}1 {è}{{\`e}}1 {ì}{{\`i}}1 {ò}{{\`o}}1 {ù}{{\`u}}1
  {À}{{\`A}}1 {È}{{\'E}}1 {Ì}{{\`I}}1 {Ò}{{\`O}}1 {Ù}{{\`U}}1
  {ä}{{\"a}}1 {ë}{{\"e}}1 {ï}{{\"i}}1 {ö}{{\"o}}1 {ü}{{\"u}}1
  {Ä}{{\"A}}1 {Ë}{{\"E}}1 {Ï}{{\"I}}1 {Ö}{{\"O}}1 {Ü}{{\"U}}1
  {â}{{\^a}}1 {ê}{{\^e}}1 {î}{{\^i}}1 {ô}{{\^o}}1 {û}{{\^u}}1
  {Â}{{\^A}}1 {Ê}{{\^E}}1 {Î}{{\^I}}1 {Ô}{{\^O}}1 {Û}{{\^U}}1
  {Ã}{{\~A}}1 {ã}{{\~a}}1 {Õ}{{\~O}}1 {õ}{{\~o}}1
  {œ}{{\oe}}1 {Œ}{{\OE}}1 {æ}{{\ae}}1 {Æ}{{\AE}}1 {ß}{{\ss}}1
  {ű}{{\H{u}}}1 {Ű}{{\H{U}}}1 {ő}{{\H{o}}}1 {Ő}{{\H{O}}}1
  {ç}{{\c c}}1 {Ç}{{\c C}}1 {ø}{{\o}}1 {å}{{\r a}}1 {Å}{{\r A}}1
  {€}{{\euro}}1 {£}{{\pounds}}1 {«}{{\guillemotleft}}1
  {»}{{\guillemotright}}1 {ñ}{{\~n}}1 {Ñ}{{\~N}}1 {¿}{{?`}}1
}

\setminted{
  fontsize=\footnotesize,
  style=xcode,
  frame=single,
  framesep=3\fboxsep,
  labelposition=topline,
}

\title{Análise e Projeto de Algoritmos}
\author{Thiago Cavalcante  -- thiago.cavalcante@penedo.ufal.br}
\institute{Universidade Federal de Alagoas -- UFAL \\ Campus Arapiraca \\ Unidade de Ensino de Penedo}
\titlegraphic{\hfill\includegraphics[height=1.5cm]{images/brasao-ufal.eps}}


\subtitle{Aula 3}
\date{08 de novembro de 2019}

\begin{document}

\maketitle

\begin{frame}{Estruturas de Dados}
  \huge
  \textbf{Tipos abstratos de dados}
  \vfill
  \begin{itemize}
    \item Contêineres (filas, pilhas)
    \item Dicionários
    \item Filas de prioridade
  \end{itemize}
\end{frame}

\begin{frame}{Estruturas de Dados}
  \huge
  Tipos abstratos de dados são implementados com \textbf{estruturas de dados}

\end{frame}

\begin{frame}{Estruturas de Dados}
  \huge
  \textbf{Estruturas contíguas $\times$ ligadas}
  \vfill
  \large
  \begin{columns}
    \begin{column}{0.49\textwidth}
      \begin{itemize}
        \item \textbf{Único} pedaço de memória
        \item Vetores (array 1D)
        \item Matrizes (array 2D)
        \item Heaps
        \item Tabelas de dispersão (\textit{hash tables})
      \end{itemize}

    \end{column}
    \begin{column}{0.49\textwidth}
      \begin{itemize}
        \item \textbf{Vários} pedaços de memória conectados por \textbf{ponteiros}
        \item Listas
        \item Árvores
        \item Listas de adjacência (grafos)
      \end{itemize}

    \end{column}
  \end{columns}
\end{frame}

\begin{frame}{Estruturas de Dados}
  \huge
  \textbf{Arrays}
  \vfill
  \LARGE
  Estruturas de \textbf{dados com tamanho fixo} onde cada elemento pode ser acessado de forma \textbf{eficicente} pelo seu \textbf{índice} (ou \textbf{endereço})

\end{frame}

\begin{frame}{Estruturas de Dados}
  \huge
  \textbf{Vantagens de um array}
  \vfill
  \Large
  \begin{itemize}
    \item Acesso aos elementos é $O(1)$, dado o índice
    \item Eficiência de espaço: composto puramente de dados com o mesmo tamanho
    \item Continuidade física dos dados na memória
  \end{itemize}

\end{frame}

\begin{frame}{Estruturas de Dados}
  \huge
  \textbf{Principal desvantagem de um array: tamanho fixo}
  \vfill
  \Large
  Possíveis soluções
  \begin{itemize}
    \item Alocar um array grande o suficiente (possível desperdício de memória)
    \item Arrays \textbf{dinâmicos} (ex.)
  \end{itemize}

\end{frame}

\begin{frame}{Estruturas de Dados}
  \huge
  \textbf{Ponteiros e estruturas ligadas}
  \vfill
  \Large
  \begin{itemize}
    \item Representam \textbf{endereços} na memória
    \item Oferecem maior \textbf{flexibilidade}
    \item Em C: operadores \textbf{\texttt{*}} e \textbf{\&}, valor \textbf{\texttt{NULL}}
  \end{itemize}

\end{frame}

\begin{frame}[fragile]{Estruturas de Dados}
  \huge
  \textbf{Implementação de uma lista}
  \vfill
  \large
  \begin{lstlisting}[language=C]
typedef struct lista {
  tipo dado; // -> dados
  Lista *prox; // -> ponteiro para o
} Lista;       //    proximo elemento
  \end{lstlisting}
\end{frame}

\begin{frame}{Estruturas de Dados}
  \huge
  \textbf{Características da lista ligada}
  \vfill
  \large
  \begin{itemize}
    \item Cada elemento armazena um ou mais dados
    \item Cada elemento possui também um \textbf{ponteiro para o próximo elemento} (ocupa mais espaço na memória)
    \item É necessário sempre o \textbf{ponteiro para o primeiro elemento} da lista
    \item Estrutura ligada mais simples
    \item Pode ser \textbf{duplamente} encadeada (ponteiros para o anterior e o próximo, ocupa mais espaço)
  \end{itemize}

\end{frame}

\begin{frame}[fragile]{Estruturas de Dados}
  \huge
  \textbf{Busca em uma lista}
  \vfill
  \large
  \begin{lstlisting}[language=C]
Lista *busca_lista(Lista *l, tipo x) {
  if (l == NULL) // lista vazia
    return(NULL);

  if (l->dado == x) // dado encontrado
    return(l);
  else // dado nao encontrado
    return(busca_lista(l->prox, x));
}
  \end{lstlisting}
\end{frame}

\begin{frame}[fragile]{Estruturas de Dados}
  \huge
  \textbf{Inserção em uma lista}
  \vfill
  \large
  \begin{lstlisting}[language=C]
void inserir_lista(Lista **l, tipo x) {
  Lista *p; // ponteiro auxiliar

  p = malloc(sizeof(Lista));
  p->dado = x;
  p->prox = *l;
  *l = p;
}
  \end{lstlisting}
\end{frame}

\begin{frame}[fragile]{Estruturas de Dados}
  \huge
  \textbf{Remoção em uma lista}
  \vfill
  \large
  \begin{lstlisting}[language=C]
Lista *anterior(Lista *l, tipo x) {
  if ((l == NULL) || (l->prox == NULL)) {
    printf("Nao encontrado");
    return(NULL);
  }
  if ((l->prox)->dado == x)
    return(l);
  else
    return(anterior(l->prox, x));
}
  \end{lstlisting}
\end{frame}

\begin{frame}[fragile]{Estruturas de Dados}
  \huge
  \textbf{Remoção em uma lista}
  \vfill
  \large
  \begin{lstlisting}[language=C]
void remover_lista(Lista **l, tipo x) {
  Lista *p, *pred;
  p = busca_lista(*l,x);
  if (p != NULL) {
    pred = anterior(*l,x);
    if (pred == NULL)
      *l = p->prox;
    else
      pred->prox = p->prox;
    free(p);
  }
}
  \end{lstlisting}
\end{frame}

\begin{frame}{Estruturas de Dados}
  \huge
  \textbf{Vantagens das listas}
  \vfill
  \large
  \begin{itemize}
    \item Uma lista nunca vai "estourar" a memória, a não ser que a memória do PC acabe
    \item Inserção e remoção são mais simples
    \item Gerenciamento de ponteiros é mais simples em larga escala
  \end{itemize}
\end{frame}

\begin{frame}{Estruturas de Dados}
  \huge
  \textbf{Desvantagens das listas}
  \vfill
  \large
  \begin{itemize}
    \item Precisam de mais espaço para os ponteiros
    \item Acesso aleatório não é eficiente
    \item Memória espalhada diminui a performance
  \end{itemize}
\end{frame}


\begin{frame}{Estruturas de Dados}
  \huge
  \textbf{Listas e arrays são estruturas \alert{recursivas}}
\end{frame}

\begin{frame}{Estruturas de Dados}
  \huge
  \textbf{Pilhas (stacks)}
  \vfill
  \Large
  \begin{itemize}
    \item \textit{last-in, first-out} (LIFO)
    \item Simples de implementar
    \item Eficientes
    \item Ordem não importa
    \item Operações: \texttt{push(p, x)} e \texttt{pop(p)}
    \item Aparecem em algoritmos recursivos
  \end{itemize}
\end{frame}

\begin{frame}{Estruturas de Dados}
  \huge
  \textbf{Filas (queues)}
  \vfill
  \Large
  \begin{itemize}
    \item \textit{first-in, first-out} (FIFO)
    \item Complicadas de implementar
    \item Minimizam o tempo máximo de espera
    \item Ordem importa
    \item Operações: \texttt{enqueue(f, x)} e \texttt{dequeue(p)}
  \end{itemize}
\end{frame}

\begin{frame}{Estruturas de Dados}
  \huge
  \textbf{Pilhas e filas podem ser implementadas tanto com arrays quanto com listas}
\end{frame}

\begin{frame}{Estruturas de Dados}
  \huge
  \textbf{Dicionários}
  \vfill
  \Large
  Armazenam um conjunto de \textbf{dados}, indexados por chaves.
  \vfill
  \large
  \texttt{\{"chave1": dado1, "chave2": dado2, $\cdots$ \}}

\end{frame}

\begin{frame}{Estruturas de Dados}
  \huge
  \textbf{Operações com dicionários}
  \vfill
  \Large
  \begin{itemize}
    \item \texttt{Busca(D, c)}
    \item \texttt{Inserção(D, x)}
    \item \texttt{Remoção(D, x)}
    \item \texttt{Máx(D) / Mín(D)}
    \item \texttt{Anterior(D, k) / Próximo(D, k)}
  \end{itemize}
  \vfill
  \large
  Exemplo: lista de elementos únicos
\end{frame}

\begin{frame}{Estruturas de Dados}
  \huge
  \textbf{Exercício: implementação de dicionário com array}
  \vfill
  Array ordenado $\times$ não ordenado
\end{frame}

\begin{frame}{Estruturas de Dados}
  \begin{center}
    \large
    \begin{tabular}{@{}lll@{}}
      \toprule
      \textbf{Operação} & \textbf{Array não ordenado} & \textbf{Array ordenado} \\
      \midrule
      Busca & \onslide<2->{$O(n)$} & \onslide<9->{$O(\log n)$} \\
      Inserção & \onslide<3->{$O(1)$} & \onslide<10->{$O(n)$} \\
      Remoção & \onslide<4->{$O(1)*$} & \onslide<11->{$O(n)$} \\
      Próximo & \onslide<5->{$O(n)$} & \onslide<12->{$O(1)$} \\
      Anterior & \onslide<6->{$O(n)$} & \onslide<13->{$O(1)$} \\
      Mín & \onslide<7->{$O(n)$} & \onslide<14->{$O(1)$} \\
      Máx & \onslide<8->{$O(n)$} & \onslide<15->{$O(1)$} \\
      \bottomrule
    \end{tabular}
  \end{center}
  \vfill
  \large
  Assumindo que o número \textbf{$n$} de elementos no array é dado
\end{frame}

\begin{frame}{Estruturas de Dados}
  \huge
  \textbf{Exercício: implementação de dicionário com lista}
  \vfill
  \Large
  Lista simplesmente ligada ordenada $\times$ lista simplesmente ligada não ordenada $\times$ lista duplamente ligada ordenada $\times$ lista duplamente ligada não ordenada
\end{frame}

\begin{frame}{Estruturas de Dados}
  \begin{center}
    \begin{tabular}{@{}lllll@{}}
      \toprule
      \textbf{Operação} & \textbf{Simp. não ord.} & \textbf{Dup. não ord.} & \textbf{Simp. ord.} & \textbf{Dup. ord.} \\
      \midrule
      Busca & \onslide<2->{$O(n)$} & \onslide<9->{$O(n)$} & \onslide<16->{$O(n)$} & \onslide<23->{$O(n)$} \\
      Inserção & \onslide<3->{$O(1)$} & \onslide<10->{$O(1)$} & \onslide<17->{$O(n)$} & \onslide<24->{$O(n)$} \\
      Remoção & \onslide<4->{$O(n)*$} & \onslide<11->{$O(1)$} & \onslide<18->{$O(n)*$} & \onslide<25->{$O(1)$} \\
      Próximo & \onslide<5->{$O(n)$} & \onslide<12->{$O(n)$} & \onslide<19->{$O(1)$} & \onslide<26->{$O(1)$} \\
      Anterior & \onslide<6->{$O(n)$} & \onslide<13->{$O(n)$} & \onslide<20->{$O(n)*$} & \onslide<27->{$O(1)$} \\
      Mín & \onslide<7->{$O(n)$} & \onslide<14->{$O(n)$} & \onslide<21->{$O(1)$} & \onslide<28->{$O(1)$} \\
      Máx & \onslide<8->{$O(n)$} & \onslide<15->{$O(n)$} & \onslide<22->{$O(1)*$} & \onslide<29->{$O(1)$} \\
      \bottomrule
    \end{tabular}
  \end{center}
\end{frame}

\begin{frame}{Estruturas de Dados}
  \huge
  \textbf{Árvores binárias de busca}
  \vfill
  \large
  \begin{itemize}
    \item Até agora: busca rápida ou atualização flexível, não ambas
    \item Busca binária requer acesso aos nós medianos \textbf{acima} e \textbf{abaixo} do que se procura
    \item Árvore binária de busca é uma "lista ligada" com dois ponteiros por elemento
  \end{itemize}

\end{frame}

\begin{frame}{Estruturas de Dados}
  \huge
  \textbf{Definição recursiva de árvore binária}
  \vfill
  \Large
  \begin{enumerate}
    \item Vazia, ou
    \item Elemento chamado de \textbf{raiz} com duas outras árvores binárias, chamadas \textbf{subárvores} da \textbf{esquerda} e da \textbf{direita}
  \end{enumerate}
  \vfill
  \large
  A ordem dos elementos importa
\end{frame}

\begin{frame}{Estruturas de Dados}
  \huge
  \textbf{Em uma árvore binária \alert{de busca}, todos os nós na \alert{subárvore da esquerda são menores} que a raiz e todos os nós na \alert{subárvore da direita são maiores} que a raiz}
  \vfill
  \large
  Com uma árvore binária de \textbf{$n$} elementos e um conjunto de \textbf{$n$} valores, só existe \textbf{uma maneira} de criar uma árvore binária de busca
\end{frame}

\begin{frame}[fragile]{Estruturas de Dados}
  \huge
  \textbf{Implementação de uma árvore}
  \vfill
  \large
  \begin{lstlisting}[language=C]
typedef struct arvore {
  tipo dado;
  struct Arvore *pai; // opcional
  struct Arvore *esquerda;
  struct Arvore *direita;
} Arvore;
  \end{lstlisting}
\end{frame}

\begin{frame}{Estruturas de Dados}
  \huge
  \textbf{Operações básicas}
  \vfill
  \LARGE
  \begin{itemize}
    \item Busca
    \item Travessia
    \item Inserção
    \item Remoção
  \end{itemize}
\end{frame}

\begin{frame}[fragile]{Estruturas de Dados}
  \huge
  \textbf{Busca em uma árvore}
  \vfill
  \large
  \begin{lstlisting}[language=C]
Arvore *busca_arvore(Arvore *l, tipo x) {
  if (l == NULL) return(NULL);
  if (l->dado == x) return(l);
  if (x < l->dado)
    return(busca_arvore(l->esquerda, x));
  else
    return(busca_arvore(l->direita, x));
}
  \end{lstlisting}
  \vfill
  \large
  Roda em tempo \textbf{$O(h)$}, onde \textbf{$h$} é a altura da árvore
\end{frame}

\begin{frame}[fragile]{Estruturas de Dados}
  \huge
  \textbf{Mínimo/máximo em uma árvore}
  \vfill
  \large
  \begin{lstlisting}[language=C]
Arvore *minimo(Arvore *t) {
  Arvore *min; // ponteiro auxiliar
  if (t == NULL) return(NULL);
  min = t;
  while (min->esquerda != NULL)
    min = min->esquerda;
  return(min);
}
  \end{lstlisting}
\end{frame}

\begin{frame}[fragile]{Estruturas de Dados}
  \huge
  \textbf{Atravessando uma árvore}
  \vfill
  \large
  \begin{lstlisting}[language=C]
void atravessar_arvore(Arvore *l) {
  if (l != NULL) {
    atravessar_arvore(l->esquerda);
    processar_elemento(l->dado);
    atravessar_arvore(l->direita);
  }
}
  \end{lstlisting}
  \vfill
  \large
  A travessia ocorre em \textbf{ordem ascendente} dos dados, em um tempo \textbf{$O(n)$} (cada um dos \textbf{$n$} elementos é processado uma vez)
\end{frame}

\begin{frame}[fragile]{Estruturas de Dados}
  \huge
  \textbf{Inserção em uma árvore}
  \vfill
  \LARGE
  Realizada no ponto onde a busca pelo elemento que se deseja inserir falha (encontra um elemento \texttt{\textbf{NULL}})
  \vfill
  \large
  Roda em tempo \textbf{$O(h)$}, onde \textbf{$h$} é a altura da árvore
\end{frame}

\begin{frame}[fragile]{Estruturas de Dados}
  \vfill
  \begin{lstlisting}[language=C]
insere_arvore(Arvore **l, tipo x, Arvore *pai) {
  Arvore *p; // ponteiro auxiliar
  if (*l == NULL) {
    p = malloc(sizeof(Arvore)); // alocacao
    p->dado = x;
    p->esquerda = p->direita = NULL;
    p->pai = pai;
    *l = p; // conectando a arvore
    return;
  }
  if (x < (*l)->dado)
    insere_arvore(&((*l)->esquerda), x, *l);
  else
    insere_arvore(&((*l)->direita), x, *l);
}
  \end{lstlisting}
\end{frame}

\begin{frame}{Estruturas de Dados}
  \huge
  \textbf{Remoção em uma árvore}
  \vfill
  \LARGE
  Três possibilidades
  \begin{enumerate}
    \item Elemento não tem filhos
    \item Elemento tem um filho
    \item Elemento tem dois filhos (complexo)
  \end{enumerate}
  \vfill
  \large
  Roda em tempo \textbf{$O(h)$}, onde \textbf{$h$} é a altura da árvore
\end{frame}

\begin{frame}{Estruturas de Dados}
  \huge
  \textbf{Pior caso}
  \vfill
  Altura \textbf{$h$} é igual ao número \textbf{$n$} de elementos, árvore vira uma lista encadeada e todos os algoritmos rodam no tempo \textbf{$O(n)$}
\end{frame}

\begin{frame}{Estruturas de Dados}
  \huge
  \textbf{Melhor caso}
  \vfill
  A árvore é \textbf{perfeitamente balanceada} (dois filhos por elemento), a altura é igual a \textbf{$\lceil \log_2 n \rceil$} e todos os algoritmos rodam no tempo \textbf{$O(\log_2 n)$}
\end{frame}

\begin{frame}{Estruturas de Dados}
  \huge
  \textbf{Na \alert{média}, existe uma \alert{alta probabilidade} de a árvore ter altura \textbf{$O(\log_2 n)$}}
\end{frame}

\begin{frame}{Estruturas de Dados}
  \huge
  \textbf{Exercício: lendo $n$ números e imprimindo na ordem com um dicionário balanceado}
  \vfill
  \large
  \begin{enumerate}
    \item Usando \texttt{inserção} e \texttt{travessia}
    \item Usando \texttt{inserção}, \texttt{mínimo} e \texttt{próximo}
    \item Usando \texttt{inserção}, \texttt{mínimo} e \texttt{remoção}
  \end{enumerate}
\end{frame}
\end{document}
